\leavevmode\hspace{-1cm}
\fontsize{14}{16}\selectfont Востряков \normalsize смотрит в окно на улицу:
Смотрю в окно и вижу снег.
 
Картина зимняя давно душе знакома.
 
Какой-то глупый человек
 
Стоит в подъезде противоположного дома.
 
Он держит пачку книг под мышкой,
 
Он курит трубку с медной крышкой.
 
Теперь он быстрыми шагами
 
Дорогу переходит вдруг,
 
Вот он исчез в оконной раме.
 
\begin{center}
(Стук в дверь).
\end{center}
 
Теперь я слышу в двери стук.
 
Кто там?
 
\leavevmode\hspace{-1cm}
\fontsize{14}{16}\selectfont Голос за дверью: \normalsize
 
Откройте. Телеграмма.
 
\leavevmode\hspace{-1cm}
\fontsize{14}{16}\selectfont Востряков \normalsize
 
Врёт. Чувствую, что это ложь.
 
И вовсе там не телеграмма.
 
Я сердцем чую острый нож.
 
Открыть иль не открыть?
 
\leavevmode\hspace{-1cm}
\fontsize{14}{16}\selectfont Голос за дверью: \normalsize
 
Откройте!
 
Чего вы медлите?
 
\leavevmode\hspace{-1cm}
\fontsize{14}{16}\selectfont Востряков \normalsize
 
Постойте!
 
Вы суньте мне под дверь посланье.
 
Замок поломан. До свиданья.
 
\leavevmode\hspace{-1cm}
\fontsize{14}{16}\selectfont Голос за дверью: \normalsize
 
Вам нужно в книге расписаться.
 
Откройте мне скорее дверь.
 
Меня вам нечего бояться,
 
Скорей откройте. Я не зверь.
 
\leavevmode\hspace{-1cm}
\fontsize{14}{16}\selectfont Востряков \normalsize (приоткрывая дверь):
 
Войдите. Где вы? Что такое?
 
\begin{center} 
(Смотрит за дверь).
\end{center}
 
Куда же он пропал? Он не мог далеко уйти. Спрятаться тут негде. Куда же он делся? 
 
Улица совсем пустая. Боже мой! И на снегу нет следов! Значит, никто к моей двери не подходил. Кто же стучал? Кто говорил со мной через дверь?
 
\begin{center}
(Закрывает дверь).
\end{center}
 
\begin{flushright}
1937 -- 1938 гг.
\end{flushright}